\chapter{Scenario Projection} \label{Chap6}

Scenario projection is structurally different from every chapter preceding it in this
dissertation. Chapters~\ref{Chap4} and~\ref{Chap5} are both checkable against real observations --
a gap-filled or forecast value can always, in principle, be compared against a genuine held-out
FCH4 measurement. Scenario projection has no such anchor: there is no ground truth for
CH\textsubscript{4} flux under a livestock density, grazing pattern, or climate trajectory that has
never actually occurred at NWFP, and there never will be, regardless of how the underlying model is
built. Every result in this chapter should be read with that structural limitation in mind; it is
returned to explicitly in Section 6.6.

\section{Outline of Scenario Projection Task}

Scenario projection extends the forecasting architecture of Chapter~\ref{Chap5} to produce
CH\textsubscript{4} projections under conditions that depart from NWFP's observed historical
record -- altered livestock stocking density, shifted grazing timing, adjusted fertiliser
application, and future climate trajectories drawn from CMIP6 projections -- rather than forecasting
forward from real, currently-observed conditions. This dissertation's scenario architecture asks a
single central "what-if" question along several axes: how does projected CH\textsubscript{4}
respond to a change in livestock management or climate, holding everything else at its
historically-observed pattern.

\section{Interpretability and Feature Importance}

The scenario architecture developed in this chapter is deliberately built around livestock density
as its central management lever, not as an unmotivated design choice. This is grounded in the
convergent interpretability evidence already established across gap-filling and forecasting
(Chapter~\ref{Chap5}): livestock density is the single feature ranked most important by every
method capable of producing an attribution on this project's model roster -- native tree-model
importance and SHAP (Random Forest, XGBoost, LightGBM), SARIMAX coefficients, and TabPFN permutation
importance -- at every tower except Tower 2, and species-disaggregated permutation importance on
this dissertation's forecasting champion further isolates that dominance to the cattle component
specifically (Chapter~\ref{Chap5}). A dedicated permutation-importance pass on TabICLv2's own
one-shot output, specific to the simplified architecture this chapter uses (Section 6.4.3), had not
yet been executed as of this writing -- a natural, low-cost follow-up flagged here rather than
assumed. The livestock-centric, species-disaggregated design of this chapter's scenario architecture
is accordingly grounded in convergent evidence from a different model and a different chapter's
work, not a fresh finding specific to this exact configuration.

\section{Simulated Climate Dataset}

Future climate projections are drawn from Semenov et al.'s large-ensemble dataset of transient
CMIP6-based climate scenarios for Great Britain \cite{semenov2025transient}, providing daily
weather series for the North Wyke location under multiple Shared Socioeconomic Pathways (SSPs) and
multiple Global Climate Models (GCMs), each contributing several stochastic realisations. This
dissertation uses SSP2-4.5 and SSP5-8.5, spanning five GCMs (ACCESS-ESM1-5, CNRM-CM6-1,
HadGEM3-GC31-LL, MPI-ESM1-2-LR, MRI-ESM2-0), transient from 2020 through to a 2050 projection
horizon.

\section{Data Pre-processing and Model Adjustment}

\subsection{Standardisation of New Data with the NWFP Dataset}

Two preprocessing steps standardise the CMIP6 series against NWFP's own historical baseline before
any scenario inference. First, of the roughly eleven driver variables this dissertation's
forecasting feature set uses, four are derived directly or near-directly from the CMIP6 series
itself (minimum, maximum, and mean air temperature; precipitation; and shortwave radiation, the
latter requiring a verified unit conversion from CMIP6's MJ/m\textsuperscript{2}/day to this
project's W/m\textsuperscript{2} convention); the remaining variables, for which CMIP6 provides no
equivalent projection, are resampled from NWFP's own historical day-of-year climatology. Second, a
direct comparison of CMIP6's own historical-period values against NWFP's real observed climate
found a genuine, non-trivial mismatch rather than a close match, in three distinct directions at
once: raw simulated air temperature runs roughly 2--3.5$^\circ$C too cool, raw simulated
precipitation is roughly 4.3 times too wet, and raw simulated incoming shortwave radiation is
roughly 8\% too low, each relative to NWFP's own real historical record. This motivated a
delta/anomaly-style bias correction, applied per GCM: an additive offset for minimum and maximum
air temperature, and a multiplicative ratio for shortwave radiation and precipitation, each
calibrated against the difference between CMIP6's own historical-period statistics and NWFP's real
observed record over the same period. The five correction factors used in this dissertation are
summarised below.

\begin{table}[h]
\centering
\resizebox{\textwidth}{!}{%
\begin{tabular}{lcccc}
\toprule
GCM & Offset TA\_min ($^\circ$C) & Offset TA\_max ($^\circ$C) & Ratio SWIN & Ratio Precipitation \\
\midrule
ACCESS-ESM1-5 & 1.81 & 2.60 & 1.034 & 0.225 \\
CNRM-CM6-1 & 2.26 & 3.09 & 1.083 & 0.229 \\
HadGEM3-GC31-LL & 1.97 & 2.79 & 1.044 & 0.236 \\
MPI-ESM1-2-LR & 2.71 & 3.56 & 1.060 & 0.226 \\
MRI-ESM2-0 & 2.29 & 3.03 & 1.043 & 0.240 \\
\bottomrule
\end{tabular}%
}
\caption{Per-GCM bias-correction factors applied to CMIP6 driver variables against NWFP's own
historical baseline.}
\end{table}

Every result in Section 6.5 uses these bias-corrected drivers rather than CMIP6's raw output. This
is a materially different, and preferred, standardisation to the coarser, output-level bias
correction used earlier in this project's scenario-analysis development (correcting a model's
predicted output to match the real historical mean, rather than correcting its input drivers) --
correcting the climate inputs themselves is the more defensible fix where the underlying mismatch is
demonstrably in the input data, rather than in the model's own extrapolation behaviour.

% \subsection{Why a Simplified, One-Shot Architecture}

% Before selecting a scenario architecture, the full forecasting roster of Chapter~\ref{Chap5} was
% rerun under a reduced-driver setting -- restricted to only the variables a genuine future climate
% scenario can actually supply, rather than the full enriched feature set Chapter~\ref{Chap5}'s
% champion has access to when forecasting from real, currently-observed conditions:

% \begin{table}[h]
% \centering
% \footnotesize
% \resizebox{\textwidth}{!}{%
% \begin{tabular}{lcccccc}
% \toprule
% Model & $R^2$ & RMSE & MAE & MASE & WAPE & Correlation \\
% \midrule
% \textbf{TabPFN} & \textbf{$-0.122$} & \textbf{56.12} & \textbf{33.14} & \textbf{0.733} & \textbf{0.899} & \textbf{0.358} \\
% TabICLv2 & $-0.330$ & 58.05 & 34.95 & 0.782 & 0.988 & 0.255 \\
% TFT & $-0.363$ & 56.59 & 35.62 & 0.841 & 1.045 & 0.292 \\
% LSTM & $-1.357$ & 63.22 & 41.06 & 0.956 & 1.268 & 0.212 \\
% SARIMAX & $-1.416$ & 65.63 & 49.54 & 1.108 & 1.516 & 0.332 \\
% XGBoost & $-1.611$ & 71.87 & 58.50 & 1.297 & 1.706 & 0.250 \\
% DLinear & $-2.068$ & 64.53 & 47.52 & 1.265 & 1.638 & 0.237 \\
% Ensemble (MASE-weighted) & $-2.101$ & 75.37 & 63.22 & 1.455 & 1.972 & 0.298 \\
% Ensemble (unweighted) & $-2.104$ & 75.31 & 63.16 & 1.454 & 1.972 & 0.299 \\
% LightGBM & $-2.682$ & 83.59 & 70.59 & 1.593 & 2.122 & 0.255 \\
% Random Forest & $-7.051$ & 99.11 & 87.77 & 2.270 & 3.121 & 0.190 \\
% \bottomrule
% \end{tabular}%
% }
% \caption{Model comparison under reduced, scenario-realistic driver availability (persistence-scaled
% MASE, distinct from Chapter~\ref{Chap5}'s climatology-scaled convention). Not the final scenario
% engine result -- a diagnostic step motivating the architecture chosen in Section 6.4.3.}
% \end{table}

% Driver availability alone reorders Chapter~\ref{Chap5}'s ranking severely: locally-fitted tree
% models collapse under reduced drivers (Random Forest falls to $R^2=-7.051$, over 28 times worse than
% its full-driver result), while the two tabular foundation models degrade far more gracefully --
% TabPFN remains the strongest model under this constraint, and TabICLv2 remains clearly usable. This
% is the direct evidence behind this chapter's model choice: a locally-fitted tree or ensemble model,
% however strong under Chapter~\ref{Chap5}'s full feature set, is not a safe choice once restricted to
% what a genuine climate scenario can actually supply.

\subsection{Generation of Simplified Scenario Model}

The scenario architecture used throughout Section 6.5 is \texttt{Direct\_TabICLv2\_solo\_trend}: a
per-tower direct tabular regression fit (\texttt{TabICLv2Regressor.fit/predict}, one fit per tower,
reused across every scenario combination for that tower) on the restricted \texttt{FX\_A\_SPECIES}
feature set plus a single linear trend covariate (\texttt{b17\_days\_since\_2010}), with no pooling
across towers and no spike-gate correction. This supersedes an earlier one-shot, no-fit
in-context configuration (\texttt{TabICLForecaster}) used in this project's initial scenario-engine
development, following B18's (Chapter~\ref{Chap5}) demonstration that direct-regression
architectures generally outperform time-series-wrapper forecasting on this task: a real-anchor
backtest (5 anchors, 3 towers, climatology-scored MASE, the same convention used throughout
Chapter~\ref{Chap5}) found every direct-regression candidate beat the prior one-shot production
architecture by a real margin, though the specific configuration mattered. A first, pooled-across-
towers candidate beat the prior architecture by 4.4\% MASE (0.726 vs.\ 0.760), but pooling assumes
one shared training cutoff across all three towers, which does not match this project's real
per-tower anchoring pipeline; a faithful, solo-per-tower adaptation without pooling recovered only
0.6\% of that margin (noise-level), until adding the trend covariate back in isolated the true
driver and recovered most of the original gain (0.738 vs.\ 0.760, \textbf{+2.79\% MASE}) -- this is
the locked-in \texttt{Direct\_TabICLv2\_solo\_trend} configuration used throughout this chapter.
The trend covariate's own training range (2018--2025, roughly 2{,}557--5{,}480 days since 2010) is
around 2.7$\times$ shorter than the horizon it is asked to extrapolate to at 2050 ($\sim$14{,}600
days) -- a genuine extrapolation risk that this dissertation checked directly rather than assumed
away: querying the fitted trend term 27--47 years past its own training maximum shows it saturates
to a bounded, plausible value ($\sim$8.3~nmol m\textsuperscript{-2} s\textsuperscript{-1}) rather
than diverging, unlike SARIMAX's known explosive extrapolation behaviour under a comparable stress
test (Chapter~\ref{Chap5}). This is a single, targeted sanity check on the trend term specifically,
not a full stress test of every scenario perturbation in Section 6.5 -- a distinction returned to in
Section 6.6.

This remains a deliberately narrower feature configuration than Chapter~\ref{Chap5}'s forecasting
champion, using the same 13 \texttt{FX\_A\_SPECIES} features as before (mean/minimum/maximum air
temperature, mean shortwave radiation, summed precipitation, day-of-year sine/cosine, growing-season
and winter indicator flags, and species-disaggregated livestock density -- combined, cattle, sheep,
and lamb) plus the trend term above, rather than the roughly 52-column enriched feature set
Chapter~\ref{Chap5}'s champion uses. It remains important to be explicit that TabICLv2 is not this
dissertation's forecasting champion -- that remains the TabPFN-based blend described in
Chapter~\ref{Chap5} (MASE=0.691) -- TabICLv2 is selected here specifically because this simplified
architecture's smaller, CMIP6-derivable feature set is a better match to what a genuine future
climate scenario can actually supply as covariates, and because the direct-regression family
degrades more gracefully than the alternatives once restricted this way (the driver-availability
ablation cited in Section 6.4.1 uses TabPFN's own reduced-driver result as the motivating evidence
for this family choice, not yet the final \texttt{FX\_A\_SPECIES} + trend configuration reported
below).

\subsection{Result of the Simplified Model}

\begin{table}[h]
\centering
\resizebox{\textwidth}{!}{%
\begin{tabular}{lc}
\toprule
Configuration & MASE (climatology) \\
\midrule
TabICLv2, BASE + ALL (Chapter~\ref{Chap5} config) & 0.7385 \\
TabICLv2, BASE + species (Chapter~\ref{Chap5} config) & 0.7405 \\
TabICLv2, FX\_A\_SPECIES, one-shot wrapper (superseded engine) & 0.7588 \\
TabICLv2, FX\_A\_SPECIES + fertiliser v2\textsuperscript{\dag} & 0.7591 \\
\textbf{Direct\_TabICLv2\_solo\_trend, FX\_A\_SPECIES + trend (this chapter's engine)} & \textbf{0.738} \\
\bottomrule
\end{tabular}%
}
\caption{Cost of simplifying to a scenario-realistic feature set, TabICLv2 only. $R^2$ is not
reported for forecasting-style comparisons anywhere in this dissertation (Section 5.1.3).
\textsuperscript{\dag}Fertiliser-variable extension tested only against the superseded one-shot
architecture, not retested against the current engine.}
\end{table}

Restricting the feature set to only what a genuine climate scenario can supply still carries a
small, quantifiable accuracy cost relative to TabICLv2's own full-feature configurations from
Chapter~\ref{Chap5} (MASE rises from 0.7405, \texttt{BASE+species}, to 0.738 under this chapter's
current engine, an increase of around 0.3\% -- a smaller gap than the superseded one-shot
architecture's own 2.5\% cost, since switching to direct regression recovers most of that
difference) -- a modest, reportable cost rather than a disqualifying one. A further
fertiliser-variable extension (\texttt{FX\_A\_SPECIES + fertiliser v2}), tested against the earlier
one-shot architecture, did not improve on it and was not adopted, consistent with this project's
general finding that fertiliser-related features add little once a richer base feature set is
already present (Section 6.5.2); it has not been retested against the current
\texttt{Direct\_TabICLv2\_solo\_trend} engine, and is not expected to change that conclusion given
the null result already found on the superseded architecture.

\section{Inference with Simulated Climate Dataset}

\subsection{Base Case}

The scenario architecture's baseline case holds livestock stocking density, grazing pattern, and
fertiliser schedule at their historically-observed day-of-year patterns, varying only the underlying
climate trajectory (SSP pathway, GCM, and realisation) drawn from the bias-corrected CMIP6 series
described above. This baseline is the reference point every induced perturbation below is measured
against.

\subsection{Induced Perturbations}

\subsubsection{Livestock Density}

Livestock density is varied along a four-point ladder, externally anchored rather than built from
an arbitrary multiplier: the historically-observed baseline (1$\times$); half that density
(0.5$\times$); a literature-informed stocking ceiling for UK grassland (drawn from stewardship
stocking-density guidance); and each tower's own historical maximum observed density. Two variants
are tested at each ladder rung: scaling all livestock species jointly, and scaling cattle alone
(motivated by Section 6.2's finding that cattle specifically drives the livestock-density signal).

\begin{table}[h]
\centering
\small
\begin{tabular}{lccc}
\toprule
Level & T2 (\% vs.\ baseline) & T4 (\% vs.\ baseline) & T9 (\% vs.\ baseline) \\
\midrule
Baseline (1$\times$) & -- & -- & -- \\
Half (0.5$\times$), all species & $-27.1\%$ & $-14.7\%$ & $-27.0\%$ \\
Half (0.5$\times$), cattle alone & $-24.6\%$ & $-14.0\%$ & $-17.6\%$ \\
Literature ceiling, all species & $+4.5\%$ & $+0.9\%$ & $-2.6\%$ \\
Literature ceiling, cattle alone & $-2.1\%$ & $+0.9\%$ & $-2.6\%$ \\
Own historical max, all species & $+98.7\%$ & $+137.8\%$ & $+182.9\%$ \\
Own historical max, cattle alone & $+57.6\%$ & $+132.5\%$ & $+187.6\%$ \\
\bottomrule
\end{tabular}
\caption{Annual-mean FCH4 change relative to the historically-observed baseline, under the
bias-corrected livestock-density ladder, this chapter's current
\texttt{Direct\_TabICLv2\_solo\_trend} scenario engine.}
\end{table}

Each tower's own historical maximum livestock density sits in a different position relative to the
literature ceiling used here, directly reflecting the differing management history described in
Chapter~\ref{Chap2}: Tower 4 and Tower 9 have both, at points in their own record, exceeded the
literature stocking ceiling used in this ladder, so their "own historical max" scenario represents a
livestock density genuinely higher than the literature-informed reference point -- consistent with
NWFP's identity as an intensively-managed research platform rather than a typical extensive UK
grassland site (Chapter~\ref{Chap2}). Tower 9's negative response at the literature-ceiling rung
($-2.6\%$) is the clearest signature of this: scaling Tower 9's livestock \emph{down} to the
literature ceiling actually reduces projected FCH4 relative to baseline, because Tower 9's real
historical stocking density already sits above that ceiling. Tower 2's own historical maximum, by
contrast, never reached this literature ceiling, reflecting its distinct land-use history (the
permanent-pasture-to-arable conversion described in Section 3.2.1). The literature-ceiling reference
value used to produce the table above is 2.5 LSU/ha, drawn from UK Countryside Stewardship
stocking-density guidance (Annex 8, non-Severely-Disadvantaged-Area land).

The scaling response is not uniformly weaker at Tower 2 across every rung of this ladder: at the
"half" and "literature-ceiling" rungs, Tower 2's proportional swing is comparable to, or larger
than, Towers 4 and 9's -- a consequence of Tower 2's much smaller absolute baseline (27.7 kg
CH4~ha\textsuperscript{-1}~yr\textsuperscript{-1}, against 72.4 and 92.2 at Towers 4 and 9
respectively, reported in full below), so the same absolute change in kg~ha\textsuperscript{-1} terms
produces a much larger percentage swing off a smaller base, not evidence of a stronger underlying
livestock response. At the "own historical maximum" rung specifically -- the most extreme, and the
only rung anchored to each tower's own real management history rather than an external reference --
Tower 2's response is decisively and consistently weaker than Towers 4 or 9's ($+98.7\%$ against
$+137.8\%$ and $+182.9\%$), consistent with Tower 2's established livestock-blindness
(Chapter~\ref{Chap4}, Section 4.5.3.1; Chapter~\ref{Chap5}, Section 5.3.4) -- a real model-behaviour
finding independently confirmed here for a fourth time, under a different model architecture and
different climate driver source than every prior confirmation, though this dissertation is explicit
that the effect is rung-dependent rather than uniform across the whole ladder.

\begin{figure}[htbp]
\centering
\includegraphics[width=\textwidth]{Figures/ch6_livestock_ladder_trajectory.png}
\caption{Monthly-mean projected FCH4 to 2050 under the livestock-density ladder (SSP2-4.5, one
representative GCM/realisation), all three towers, this chapter's current scenario engine.
Own-historical-max visibly separates from the other three rungs at every tower.}
\end{figure}

\FloatBarrier

\paragraph{Annual CH4 generation in absolute terms.}

The percentage changes above are all relative to each tower's own historical baseline, which
differs sharply in absolute magnitude across towers and can make a small absolute change look large,
or vice versa (the literature-ceiling discussion above is one direct example). The livestock-density
ladder is reported below in absolute annual CH4 generation terms instead:

\begin{table}[h]
\centering
\small
\begin{tabular}{lccc}
\toprule
Level & T2 (kg~ha\textsuperscript{-1}~yr\textsuperscript{-1}) & T4 (kg~ha\textsuperscript{-1}~yr\textsuperscript{-1}) & T9 (kg~ha\textsuperscript{-1}~yr\textsuperscript{-1}) \\
\midrule
Baseline (1$\times$) & 27.7 & 72.4 & 92.2 \\
Half (0.5$\times$), all species & 20.2 & 61.7 & 67.3 \\
Literature ceiling, all species & 29.0 & 73.0 & 89.7 \\
Own historical max, all species & 55.1 & 172.1 & 260.7 \\
\bottomrule
\end{tabular}
\caption{Projected annual CH4 generation (kg CH4 per hectare per year) under the livestock-density
ladder, SSP2-4.5, this chapter's current scenario engine.}
\end{table}

In absolute terms, Tower 9's baseline and own-historical-max generation both exceed Tower 4's, which
in turn exceeds Tower 2's, at every rung of the ladder -- the same tower ordering the percentage
table above can obscure at the smaller rungs. Own-historical-max remains the largest single lever
tested anywhere in this chapter in absolute terms too, reaching 260.7 kg~ha\textsuperscript{-1}~yr\textsuperscript{-1}
at Tower 9, nearly three times its own baseline.

\begin{figure}[htbp]
\centering
\includegraphics[width=\textwidth]{Figures/ch6_annual_ch4_livestock.png}
\caption{Projected annual CH4 generation (kg CH4 per hectare per year) under the livestock-density
ladder, SSP2-4.5, all three towers. Grazing-timing and fertiliser-schedule perturbations (Section
6.5.2) produce the analogous absolute-terms figures, not reproduced here since their percentage
effects are already small (Table above, Section 6.5.2).}
\end{figure}

\FloatBarrier

\subsubsection{Grazing Timing and Fertiliser Schedule}

\begin{table}[h]
\centering
\begin{tabular}{lccc}
\toprule
Intervention & T2 & T4 & T9 \\
\midrule
Extend grazing season by 4 weeks & $-3.8\%$ & $+19.9\%$ & $+19.9\%$ \\
Fertiliser frequency $+50\%$ & $-8.2\%$ & $-0.1\%$ & $\approx 0\%$ \\
Fertiliser rate $+50\%$ & $-1.7\%$ & $-2.0\%$ & $\approx 0\%$ \\
\bottomrule
\end{tabular}
\caption{Annual-mean FCH4 change relative to baseline under grazing-timing and fertiliser
perturbations (SSP2-4.5), this chapter's current \texttt{Direct\_TabICLv2\_solo\_trend} scenario
engine.}
\end{table}

Two further management perturbations are tested against the same bias-corrected climate driver
pipeline. Shifting the grazing season four weeks later still produces a real, monotonic increase in
projected FCH4 at Towers 4 and 9 (+19.9\% each, materially larger than this same perturbation's
+18.6\% under the earlier one-shot scenario engine, though the qualitative conclusion -- a real,
management-actionable lever -- is unchanged). Tower 2's response, however, \emph{reverses sign}
under this chapter's current engine: $-3.8\%$, against $+3.8\%$ under the earlier architecture. This
dissertation reports this reversal directly rather than smoothing over it: Tower 2's grazing response
was already the smallest-magnitude and least livestock-driven result in this table under the
earlier architecture, consistent with Tower 2's established livestock-blindness
(Chapter~\ref{Chap4}, Section 4.5.3.1; Chapter~\ref{Chap5}, Section 5.3.4), so a sign flip at this
small a magnitude is consistent with Tower 2 genuinely lacking a stable livestock-driven signal to
extend, rather than the underlying architecture change producing a materially different real-world
finding -- Towers 4 and 9's own management-actionable conclusion is unaffected either way. Fertiliser
schedule perturbations (a 50\% increase in either application frequency or per-event rate) produce a
small, still sign-inconsistent-across-towers effect, now reaching as much as $-8.2\%$ at Tower 2
(frequency) rather than the earlier architecture's $-3.5\%$ -- a larger magnitude than before, but
still consistent with this project's standing conclusion that fertiliser timing carries little
independent signal once livestock density and the richer base feature set are already present, and
still the smallest management lever tested in this chapter for Towers 4 and 9 specifically ($\leq
2.0\%$ in either direction). Separately, the SSP2-4.5 vs.\ SSP5-8.5 pathway divergence is small
relative to every management perturbation above throughout most of the projection horizon, though it
grows toward the end of the horizon and reaches its largest value at Tower 2 (approximately 2.9\%
divergence late in the horizon) -- a minor lever next to the livestock and grazing-timing effects
reported above, but not strictly negligible everywhere.

\begin{figure}[htbp]
\centering
\includegraphics[width=\textwidth]{Figures/ch6_grazing_trajectory.png}
\caption{Monthly-mean projected FCH4 to 2050 under the grazing-timing shift (SSP2-4.5, one
representative GCM/realisation), all three towers, this chapter's current scenario engine. The
shifted-timing curve sits visibly above the historical-timing curve at Towers 4 and 9, and is
slightly below it at Tower 2.}
\end{figure}

\FloatBarrier

\section{Limitations}

Three limitations apply specifically to this chapter's scenario projections, beyond the general
absence of ground truth stated in its opening framing.

First, TabICLv2's extrapolation behaviour under genuine distribution shift has not been formally
stress-tested to the standard Chapter~\ref{Chap5}'s uncertainty-quantification stress test applied
to the production ensemble -- TabICLv2 was not part of that model roster, and this chapter's
livestock-ladder sweep, while exposing the model to out-of-historical-range livestock density, is a
behavioural probe rather than a controlled extrapolation test with any ground truth to check
against. TabICLv2 belongs to the TabPFN-family of in-context tabular models, a family with a
documented structural extrapolation limitation under genuine distribution shift
(Section~\ref{sec:tabular-modelling}). One narrower piece of this concern was checked directly rather
than left as an open risk: the current engine's trend covariate (Section 6.4.3) is fitted on roughly
2018--2025 data and queried out to 2050, a genuine extrapolation in time, and a direct sanity check
found it saturates to a bounded, plausible value rather than diverging when queried decades past its
own training range. This is reassuring for the time-trend extrapolation specifically, but it is a
single, targeted check on one covariate, not a substitute for a full stress test of the model's
behaviour under livestock-density or climate-driver distribution shift -- the two forms of
extrapolation this chapter's scenarios actually rely on most heavily. This remains a limitation of
this chapter's headline result specifically, not a secondary footnote, since the livestock-density
and climate scenarios reported above are precisely the kind of distribution shift this limitation
concerns.

Second, no diagnostic available to this dissertation can detect a smooth-but-confidently-wrong
extrapolation. This is a structural gap rather than a tooling gap that could simply be built:
historical PICP (Chapter~\ref{Chap5}) measures calibration only at real, historically-anchored
points, and cannot be computed at all for a genuine future scenario, since no future ground truth
exists to check against; a narrow prediction interval width is fully compatible with a confidently
wrong extrapolation; and GCM/realisation spread (Section 6.3) measures sensitivity to the climate
input alone, not error in the underlying model's own extrapolation. One further distinction is
stated explicitly here to avoid it being misread elsewhere in this dissertation: prediction
intervals were successfully generated for over 99\% of eligible scenario rows in this chapter's
pipeline, but this is \emph{interval availability} (the pipeline ran without failure), not
\emph{empirical coverage} -- it says nothing about whether those intervals are actually well
calibrated under scenario conditions, for exactly the reason given above.

Third, the CMIP6-vs-NWFP historical mismatch identified in Section 6.4.1 is a real, separate
limitation independent of model choice -- it reflects a genuine discrepancy between the global
climate model ensemble's own historical baseline and NWFP's real local climate record, corrected for
as best this dissertation can (Section 6.4.1), but not eliminated as a source of residual
uncertainty in the corrected driver series feeding every result in Section 6.5.
